\section{The same boundary trap for the uncapped surrogate step}
\label{app:bsur-trap}

This section considers a second actual trajectory on the finite instance of
Section~\ref{app:trap}. The slope certificate, relative margin, and exact-refresh
semantics are unchanged. Only the curvature coefficient in the accepted step is
replaced by
\[
 R_2=\sum_iw_i\left(\frac{p_i^2}{q}-1\right)^2p_i^2,
 \qquad B_{\rm sur}=\left(\sqrt{R_2}
                  +\gamma\sqrt{Q/\alpha}\right)^2.
\]
Thus a cache is retained exactly when $g>0$ and $G\ge g/2$. Otherwise it is
recomputed at the current design; only a fresh $g=0$ terminates the run. An
accepted update uses
\begin{equation}\label{eq:bsur-rule}
 t=\min\{1/2,G/(4B_{\rm sur})\},\qquad w^+=w+td.
\end{equation}
The quotient is omitted if $B_{\rm sur}=0$, a case that will not occur in the
construction. There is no refresh-count cap. The old-cache quantities in this
rule are exactly
\[
 q=\sum_iw_ip_i^2,\quad d_i=w_i(p_i^2/q-1),\quad
 g=\frac{\operatorname{Var}_w(p_i^2)}q,\quad
 G=g-\frac{2\sqrt{HQ}+\zeta Q}{\alpha}.
\]
In particular, the cached fit is not silently replaced by a current fit when
evaluating the gate or the step.

\paragraph{Claim and fixed instance.}
Set $R=100$, $B=101/100$, $\delta=99/100$, and $N=2^{10^7}$. The rows
\[
 (1,0),\quad(\sqrt\delta,0),\quad
 (\sqrt B,-R),\quad(\sqrt B,1)
\]
have respective multiplicities $N^3,N^2,1,N$, and all
$m=N^3+N^2+N+1$ individual rows initially have weight $1/m$.
For the rule~\eqref{eq:bsur-rule}, every finite weight stays positive, every cache
epoch is finite, exact refreshes occur infinitely often, and
\begin{equation}\label{eq:bsur-claim}
 v(w_k)\uparrow1<101/100=v_*,\qquad
 \operatorname{eff}(w_k)\longrightarrow100/101.
\end{equation}
Here $k$ counts accepted updates. The assertion permits repeated rows and does
not assert uniform initialization on the four distinct types after deduplication.

Equal copies remain equal under the update. All sums, now including $R_2$, and
all extrema defining $\alpha,\gamma,\zeta$ aggregate exactly. The four aggregate
masses $(w_0,\omega,x_1,x_2)$ therefore start at
\[
 (w_0,\omega,x_1,x_2)=m^{-1}(N^3,N^2,1,N).
\]
The feature matrix has full column rank, and direct elimination gives
\begin{equation}\label{eq:bsur-value}
 A_w=R^2x_1+x_2>0,\qquad
 v(w)=w_0+\delta\omega+
 B(R+1)^2\frac{x_1x_2}{R^2x_1+x_2}.
\end{equation}
Evaluation at $\theta=0$ bounds every design value by $B$. The informative-pair
design $(x_1,x_2)=(1/(R+1),R/(R+1))$ attains this bound, so $v_*=B$.
The two zero-nuisance types have distinct fixed squared residuals $1,\delta$;
their positive weights ensure $g>0$ at every finite state, including refreshes.

\subsection{An actual Bsur epoch: geometry and finite exit}
\label{sec:bsur-geometry}

At a fresh cache, let $x$ denote the initially smaller informative mass and $y$
the other one, keeping these labels fixed throughout the epoch. Its type is
$\mathrm H$ if $x=x_1$ and $\mathrm L$ if $x=x_2$. Assume
\begin{equation}\label{eq:bsur-entry}
 0<\omega_s\le10^{-6},\qquad
 s_y/\omega_s=e^{-b},\qquad s_x/\omega_s=e^{-2b-\varepsilon},
 \qquad b\ge1000,\quad |\varepsilon|\le100.
\end{equation}
The epoch ends at the first state where the old cache fails its relative margin,
before taking a fresh fit there. We first prove invariance, a positive step
lower bound, and finite exit for the actual Bsur steps.

Write $\rho=s_x/s_y=e^{-b-\varepsilon}$,
$\kappa=z_x^2/z_y^2\in\{10^4,10^{-4}\}$, and
\[
 \mathcal A_{\mathrm H}=1030301/100,\qquad
 \mathcal A_{\mathrm L}=1030301/10^6.
\]
The fresh scalar normal equation gives the fixed cached scores
\begin{equation}\label{eq:bsur-scores}
 P=p_x^2=\frac{\mathcal A_\sigma}{(1+\kappa\rho)^2},\qquad
 S=p_y^2=\kappa\rho^2P,\qquad
 A_s=z_y^2s_y(1+\kappa\rho).
\end{equation}
In particular,
\[
 1.02<P<11000,\qquad 0<S<10^{-20},\qquad
 1\le1+\kappa\rho\le2,\qquad S<\delta<1<P.
\]
These are fresh-fit identities; they do not assert equality with the trajectory
using $K$.

\paragraph{Conditional regional bounds.}
First suppose a current state belongs to the region
\begin{equation}\label{eq:bsur-region}
 \omega\le10^{-6},\qquad x/\omega\le10^{-20},\qquad y/\omega\le10^{-20}.
\end{equation}
An induction below proves this condition through the first failed endpoint.
Since $w_0=1-\omega-x-y$,
\[
 q=1-.01\omega+(P-1)x+(S-1)y,
 \qquad 1-.011\omega\le q\le1-.009\omega,
 \qquad .999<q<1.
\]
The auxiliary term in $qg=\sum_iw_i(p_i^2-q)^2$ gives
$g\ge10^{-5}\omega$. For an upper bound, its contribution to $g$ is below
$.000101\omega$, the base contribution is at most
$.011^2\omega^2/.999$, and the other contributions total at most
$(11000^2+1)10^{-20}\omega/.999$. Thus
\begin{equation}\label{eq:bsur-g-bounds}
 10^{-5}\omega\le g<.000102\omega<10^{-3}\omega.
\end{equation}
The new curvature statistic is exactly
\[
 R_2=\frac{w_0(1-q)^2+\delta\omega(\delta-q)^2
                 +Px(P-q)^2+Sy(S-q)^2}{q^2}.
\]
Put $U=R_2/g$. The auxiliary contribution gives
$R_2>\delta(.009)^2\omega=.00008019\omega$;
since $.00008019>(3/4)(.000102)$, we have $U>3/4$.
For the upper bound use the exact centered-moment identity
\[
 q^2(R_2-g)=\sum_iw_i(\lambda_i-q)^3,
 \qquad \lambda=(1,\delta,P,S).
\]
The auxiliary contribution is less than $-(.009)^3\omega$, the base contribution
is at most $.011^3\omega^3$, the $x$ contribution is at most
$11000^3\,10^{-20}\omega$, and the $y$ contribution is negative. Hence
\[
 q^2(R_2-g)<\omega[-.009^3+.011^3\,10^{-12}+11000^3\,10^{-20}]<0.
\]
We have proved the step-independent regional bounds
\begin{equation}\label{eq:bsur-ratio-region}
 3/4<U<1,
\end{equation}
and $R_2>0$.

For $0<t\le1/2$, let $F_V=1+t(V/q-1)$ denote a row multiplier. The same
regional bounds give
\begin{equation}\label{eq:bsur-multipliers}
 \frac{F_P}{F_\delta}<10^4,\quad
 \frac{F_S}{F_\delta}<1,\quad
 \frac{F_P^2}{F_\delta F_S}<10^{10},\quad
 F_\delta\le1-.009t.
\end{equation}
For example, $F_P<5600$, $F_\delta\ge.995$, and $F_S\ge1/2$ suffice for the
two upper ratio bounds.

\paragraph{The actual fresh step and informative growth.}
At the fresh origin, $Q=0$, $G=g$, and $B_{\rm sur}=R_2$. Therefore
\begin{equation}\label{eq:bsur-fresh-step}
 t_{\rm fresh}=1/(4U)\in(1/4,1/3),\qquad
 \omega^+<.99775\omega.
\end{equation}
This step is not a half-step. After it,
\[
 J:=\frac{x/s_x}{y/s_y}
 =\frac{F_P}{F_S}
 =1+\frac{t(P-S)}{q(1-t)+tS}>4/3.
\]
Indeed, the numerator of $J-1$ exceeds $(1/4)(1.02-10^{-20})$ and its
denominator is below $3/4+(1/3)10^{-20}$. Three times that lower numerator
exceeds the upper denominator. Throughout the remaining epoch, $J$ increases.
In a low-type fresh epoch $J<1+(1/2)\mathcal A_{\rm L}/.999<1.52$;
thus the earlier $K$ proof's $J>2$ cannot be used for this rule.

\paragraph{The actual gate and step floor.}
The ordering of the fixed scores preserves the ordering of current-to-cache
ratios, including the base and auxiliary types. Thus, with
$\eta_x=x/s_x$ and $\eta_y=y/s_y$, we have $\alpha=\eta_y$.
Put $\mathcal C=z_xp_xs_x=-z_yp_ys_y$ and
$\Delta_x=P/q-1$, $\Delta_y=S/q-1$. At a nonfresh state,
\[
 r=\mathcal C(\eta_x-\eta_y),\qquad
 h=\mathcal C(\eta_x\Delta_x-\eta_y\Delta_y),
\]
so
\[
 \ell:=h/r=\Delta_x+\frac{\Delta_x-\Delta_y}{J-1},
 \qquad 0<\ell<4P/q<10^5.
\]
Also $\gamma<10^5$ and $\zeta=1-S/q\in[1/2,1]$. In particular
$\sqrt H=\ell\sqrt Q$ without a sign ambiguity. Direct substitution gives
\begin{equation}\label{eq:bsur-qalpha}
 \frac Q\alpha=
 \frac{\kappa P}{1+\kappa\rho}(1-J^{-1})^2\frac{x^2}{y}.
\end{equation}
Its coefficient lies between $10^{-6}$ and $10^9$: for the lower bound use
$1-J^{-1}>1/4$ and
$(10^{-4})(1.02)/(2\cdot16)>10^{-6}$.
Define $W=x^2/(\omega y)$ and $\mathcal Z=Q/(\alpha g)$. We obtain
\begin{equation}\label{eq:bsur-gate-ratios}
 10^{-3}W\le\mathcal Z\le10^{14}W,\qquad
 1-G/g=(2\ell+\zeta)\mathcal Z\le10^6\mathcal Z,
 \qquad
 \frac{B_{\rm sur}}g=(\sqrt U+\gamma\sqrt{\mathcal Z})^2.
\end{equation}
The last expression retains the nonzero fresh contribution $U$; it is not the
corresponding expression for $K$.
An accepted nonfresh origin has $\mathcal Z\le1$ and $W\le10^3$. By
\eqref{eq:bsur-ratio-region}, its curvature ratio is below
$(1+10^5)^2<4\cdot10^{10}$. Thus its actual selected step satisfies
\begin{equation}\label{eq:bsur-step-floor}
 t\ge\min\{1/2,(32\cdot10^{10})^{-1}\}>10^{-12}.
\end{equation}
The fresh step obeys the same lower bound by~\eqref{eq:bsur-fresh-step}.

\paragraph{Invariance and finite exit.}
The initial ratios in~\eqref{eq:bsur-entry} are far below $10^{-20}$.
Its first actual Bsur step preserves the region by
\eqref{eq:bsur-multipliers}. Subsequently $y/\omega$ decreases and remains at
most $e^{-b}$. At any accepted nonfresh origin already in the region,
\[
 x/\omega=\sqrt{W(y/\omega)}\le\sqrt{10^3}e^{-b/2}.
\]
Its next update gives $x^+/\omega^+<10^4\sqrt{10^3}e^{-b/2}<10^{-20}$;
the other ratio and $\omega$ decrease. This proves the region by induction,
including the first failed endpoint. It does not require that endpoint to pass
the gate.

For each accepted step put
\[
 a_k=\log(F_P/F_\delta)>0,\qquad
 \beta_k=\log(F_\delta/F_S)>0.
\]
Their exact increments satisfy
\begin{equation}\label{eq:bsur-log-motion}
 \log W_{k+1}-\log W_k=2a_k+\beta_k,
 \qquad \beta_k\ge\log(1+.98t_k)\ge t_k/2.
\end{equation}
The latter follows from $F_\delta-F_S>.98t_k$ and $F_S\le1$.
Every accepted step therefore increases $\log W$ by more than
$5\cdot10^{-13}$, whereas every accepted origin has $W\le10^3$, including
$W_0=e^{-3b-2\varepsilon}$ at the fresh state. The epoch must end after finitely
many accepted steps $M$; a conservative explicit bound is
\[
 M\le2+\left\lceil2\cdot10^{12}
                  (3b+200+\log10^3)\right\rceil.
\]
At the first failed endpoint, $g>0$ and $G/g<1/2$, so
\[
 1/2<10^6\mathcal Z\le10^{20}W_{\rm end}.
\]
The preceding accepted origin has $W\le10^3$, and its final multiplier for W
is below $10^{10}$. Consequently
\begin{equation}\label{eq:bsur-endpoint}
 10^{-22}<W_{\rm end}\le10^{13},\qquad
 |\log W_{\rm end}|\le100.
\end{equation}
No current fit, approximate refresh threshold, or assumed exit time was used.

\subsection{Summable error over the actual variable-step prefix}
\label{sec:bsur-prefix}

Use the reference step $t_0=1/(4\delta)=25/99$ and define
\[
 a_\sigma(t)=\log\frac{1+t(\mathcal A_\sigma-1)}{1+t(\delta-1)},
 \qquad
 \beta(t)=\log\frac{1+t(\delta-1)}{1-t}.
\]
These are reference increments; no actual step is assumed to equal $t_0$.
Write $a_\sigma=a_\sigma(t_0)$ and $\beta=\beta(t_0)$.

\paragraph{A summable bound on $U-\delta$.}
Since $qg=\sum_iw_i(\lambda_i-q)^2$,
\[
 U-\delta/q=
 \frac{\sum_iw_i(\lambda_i-\delta)(\lambda_i-q)^2}{q^2g}.
\]
The auxiliary numerator is zero. Using $q^2g\ge.999^2\,10^{-5}\omega$,
$1-q\le.011\omega$, and the score bounds yields
\begin{equation}\label{eq:bsur-u-error}
 |U-\delta|\le\omega+2\cdot10^{17}x/\omega+2\cdot10^5y/\omega.
\end{equation}
For clarity, the three coefficients before rounding upward are bounded by
\[
 \frac{\delta(.011)}{.999}
 +\frac{.01(.011)^2}{.999^2\,10^{-5}}<1,\qquad
 \frac{11000^3}{.999^2\,10^{-5}}<2\cdot10^{17},\qquad
 \frac1{.999^2\,10^{-5}}<2\cdot10^5.
\]

\paragraph{The prefix and its actual steps.}
Let $W_{\rm cut}=10^{-40}$ and let $n$ be the first index with
$W_n>W_{\rm cut}$. It exists before or at the failed endpoint: initially W is
below the cutoff, the failed endpoint is above it, and all states below the
cutoff will pass the gate. At any prefix origin $k<n$,
\[
 \mathcal Z\le10^{-26},\qquad
 \eta:=1-G/g\le10^{-20},\qquad
 \gamma\sqrt{\mathcal Z}\le10^{-8}.
\]
At the fresh origin these upper bounds hold with $\mathcal Z=\eta=0$, without
defining $h/r$ there. Put
\[
 V=B_{\rm sur}/g=(\sqrt U+\gamma\sqrt{\mathcal Z})^2.
\]
Weakening the regional upper bound to $U<3/2$ gives
$3/4<V<(5/4+10^{-8})^2<7/4$. The prefix gate passes, and the half-step cap is
inactive. Thus the actual step obeys
\begin{equation}\label{eq:bsur-prefix-step}
 t_k=\frac{1-\eta}{4V},\qquad 1/8<t_k<1/3.
\end{equation}
Its discrepancy from the reference step is bounded by
\begin{align*}
 |t_k-t_0|
 &\le\frac{\eta}{4V}+\frac{|V-\delta|}{4\delta V}\\
 &\le\frac\eta3+\frac12|U-\delta|
       +\gamma\sqrt U\sqrt{\mathcal Z}
       +\frac12\gamma^2\mathcal Z.
\end{align*}
The cross term is at most $(5/4)10^{12}\sqrt W$. The first and last terms
together are at most $(10^6/3+10^{10}/2)10^{14}W$.
Since $W\le10^{-40}$ implies $W\le10^{-20}\sqrt W$, these estimates give
\begin{equation}\label{eq:bsur-step-error}
 |t_k-t_0|\le |U-\delta|+2\cdot10^{12}\sqrt{W_k}.
\end{equation}

\paragraph{Uniform sums, without multiplying by the epoch length.}
By~\eqref{eq:bsur-log-motion} and~\eqref{eq:bsur-prefix-step}, each prefix
step raises $\log W$ by at least $1/16$. Summing backward from
$W_{n-1}\le10^{-40}$, and using the forward contraction of $y/\omega$, gives
\begin{equation}\label{eq:bsur-prefix-sums}
 \begin{aligned}
 \sum_{k<n}\sqrt{W_k}&\le33\cdot10^{-20},&
 \sum_{k<n}\omega_k&\le(8000/9)\omega_s,\\
 \sum_{k<n}y_k/\omega_k&\le17e^{-b},&
 \sum_{k<n}x_k/\omega_k&\le33\cdot10^{-20}e^{-b/2}.
 \end{aligned}
\end{equation}
Indeed $1/(1-e^{-1/32})<33$ bounds the first geometric sum. The second uses
$\omega_{k+1}\le(1-9/8000)\omega_k$, and the third uses
$1/(1-e^{-1/16})<17$. These denominator bounds follow from
$e^{-u}\le1/(1+u)$. The fourth estimate uses the identity
$x/\omega=\sqrt{W(y/\omega)}$ and $y/\omega\le e^{-b}$.

Combining~\eqref{eq:bsur-u-error}, \eqref{eq:bsur-step-error}, and
\eqref{eq:bsur-prefix-sums}, for every $b\ge1000$,
\begin{equation}\label{eq:bsur-total-step-error}
 \begin{aligned}
 \sum_{k<n}|t_k-t_0|
 &\le(8000/9)10^{-6}
 +(2\cdot10^{17})(33\cdot10^{-20})e^{-500}\\
 &\quad +(2\cdot10^5)(17)e^{-1000}
 +(2\cdot10^{12})(33\cdot10^{-20})
 <10^{-3}.
 \end{aligned}
\end{equation}
The last inequality remains strict on replacing the exponentials by the larger
$2^{-500},2^{-1000}$. Also,
\begin{equation}\label{eq:bsur-total-q-error}
 \sum_{k<n}|q_k-1|\le.011\sum_{k<n}\omega_k<10^{-5}.
\end{equation}
These bounds control the entire prefix independently of its length.

\subsection{The logarithmic epoch map and the original finite start}
\label{sec:bsur-map}

The last prefix step has $W_n\le10^{10}W_{\rm cut}=10^{-30}<1$.
Its accumulated log-W motion is at least $n/16$, while
$-\log W_0=3b+2\varepsilon$. Therefore
\[
 n<16(3b+200)\le100b.
\]
The fixed-cache residual errors in~\eqref{eq:bsur-scores} satisfy
\[
 |P-\mathcal A_\sigma|\le2.2\cdot10^8e^{-b+100},\qquad
 S\le1.1\cdot10^8e^{-2b+200}.
\]
Consequently
\begin{equation}\label{eq:bsur-total-residual-error}
 n|P-\mathcal A_\sigma|
 \le100b(2.2\cdot10^8)e^{-b+100}<1,
 \qquad nS\le100b(1.1\cdot10^8)e^{-2b+200}<1.
\end{equation}
Both upper bounds decrease for $b\ge1000$; at $b=1000$ even the larger powers
$2^{-900},2^{-1800}$ suffice. Multiplication by the prefix length is legitimate
here because that length has been bounded and the cache errors decay
exponentially in the entry depth.

The exact variable-step logarithmic increments are
\[
 \begin{aligned}
 a_k&=\log[q_k(1-t_k)+t_kP]-\log[q_k(1-t_k)+t_k\delta],\\
 \beta_k&=\log[q_k(1-t_k)+t_k\delta]-\log[q_k(1-t_k)+t_kS].
 \end{aligned}
\]
For $q\in[.999,1]$, $t\in[1/8,1/3]$, and the intervening nonnegative score
ranges, either log ratio has derivatives bounded in absolute value by 16 in t,
3 in q, and 1 in its informative score P or S. For example,
\[
 \left|\partial_t\log(q(1-t)+tU)\right|
 \le\max\{1/t,1/(1-t)\}\le8,\qquad
 0\le\partial_q\log(q(1-t)+tU)\le1/q.
\]
The score derivative is at most $t/[q(1-t)]<1$.
The same bounds hold on the comparison paths to $q=1$, $t=t_0$,
$P=\mathcal A_\sigma$, and $S=0$.
Equations~\eqref{eq:bsur-total-step-error}--\eqref{eq:bsur-total-residual-error}
therefore give
\begin{equation}\label{eq:bsur-prefix-log-error}
 \left|\sum_{k<n}a_k-na_\sigma\right|<2,\qquad
 \left|\sum_{k<n}\beta_k-n\beta\right|<2.
\end{equation}
Thus the actual varying prefix has been controlled, rather than replaced by a
fixed-step recurrence.

\paragraph{The unrestricted tail and endpoint reset.}
No step approximation is used after the cutoff. Positivity of both logarithmic
increments and~\eqref{eq:bsur-endpoint} imply
\[
 \sum_{k=n}^{M-1}(2a_k+\beta_k)
 =\log(W_{\rm end}/W_n)<53\log10<159.
\]
The tail contributes less than $79.5$ to the a sum and $159$ to the beta sum,
regardless of the number and sizes of its actual steps. With
$X=-\log(x/\omega)$ and $Y=-\log(y/\omega)$, we consequently have
\[
 X_{\rm end}=2b+\varepsilon-na_\sigma+E_X,\qquad
 Y_{\rm end}=b+n\beta+E_Y,\qquad |E_X|,|E_Y|<200.
\]
Set $b'=X_{\rm end}$ and
$\varepsilon_{\rm next}=Y_{\rm end}-2X_{\rm end}=\log W_{\rm end}$.
Eliminating n gives
\begin{equation}\label{eq:bsur-exact-map}
 \begin{aligned}
 b'={}&C_\sigma b
 +\frac{\beta}{2a_\sigma+\beta}\varepsilon
 -\frac{a_\sigma}{2a_\sigma+\beta}\varepsilon_{\rm next}\\
 &+\frac{a_\sigma}{2a_\sigma+\beta}E_Y
 +\frac{\beta}{2a_\sigma+\beta}E_X,
 \qquad C_\sigma=\frac{a_\sigma+2\beta}{2a_\sigma+\beta}.
 \end{aligned}
\end{equation}
Every error coefficient has modulus below one. Hence the actual Bsur epoch obeys
\begin{equation}\label{eq:bsur-map-bounds}
 |b'-C_\sigma b|<1000,\qquad |\varepsilon_{\rm next}|\le100.
\end{equation}
The first bound follows already with 600. The second is a new endpoint bound,
not an accumulation of the old epsilon parameters. If $b'\ge1000$, then
$2b'+\varepsilon_{\rm next}>b'$, so the informative rows really exchange
their small and large roles when the next epoch is relabeled.

\paragraph{A finite two-epoch margin.}
At $t_0=25/99$ the constants simplify to
\[
 a_{\mathrm H}=\log(1030597/395),\qquad
 a_{\mathrm L}=\log(3990301/3950000),\qquad
 \beta=\log(395/296).
\]
The following strict finite bounds suffice:
\[
 a_{\mathrm H}<7.87,\qquad a_{\mathrm L}<.01016,\qquad \beta>.2885.
\]
For the first, the positive Taylor polynomial
$\sum_{j=0}^{20}(7.87)^j/j!$ exceeds $1030597/395$ and is below $e^{7.87}$.
For the second, use $\log(1+x)\le x-x^2/2+x^3/3$ at
$x=40301/3950000$. For the third, put $z=99/691$ and use
\[
 \beta=2\operatorname{arctanh}z
 >2(z+z^3/3)=95188104/329939371>.2885.
\]
The function $C(a,\beta)=(a+2\beta)/(2a+\beta)$ decreases in a and increases
in beta. Thus
\begin{equation}\label{eq:bsur-map-expansion}
 C_{\mathrm H}C_{\mathrm L}>
 \frac{7.87+.577}{15.74+.2885}
 \frac{.01016+.577}{.02032+.2885}
 =\frac{165324684}{164997379}>1.0019,
 \qquad 1/2<C_\sigma<2.
\end{equation}
This is a smaller expansion than for K, and requires its own finite entry check.

Suppose an H-type epoch starts with $b_H\ge5\cdot10^6$. Its next parameter
satisfies $b_L>b_H/2-1000\ge2499000>1000$, so the actual role exchange is
valid and the next L-type epoch satisfies~\eqref{eq:bsur-entry}.
After that L epoch, the lower bound $C_{\mathrm L}>1/2$ first gives a new
parameter greater than $1248500$, independently verifying the second role
exchange. Combining the two maps more sharply gives
\begin{equation}\label{eq:bsur-global-induction}
 b_{H,\rm next}>1.0019b_H-3000\ge1.001b_H.
\end{equation}
The final inequality has at least 1500 of slack at $b_H=5\cdot10^6$.
The parameter error in b has therefore been charged at both epochs; the
endpoint epsilon bound is not used to suppress it.

For the unchanged individual-row uniform initialization,
\[
 x_1/\omega=N^{-2},\qquad x_2/\omega=N^{-1},\qquad
 b_0=10^7\log2>5\cdot10^6,\qquad
 \varepsilon_0=0,\qquad \omega_0<1/N<10^{-6}.
\]
It enters the H-type induction exactly. All subsequent H parameters increase
geometrically and every intervening L parameter remains above 1000.
Equations~\eqref{eq:bsur-entry}--\eqref{eq:bsur-map-bounds} can therefore be
iterated indefinitely, with the original finite N and no modified initialization.
Each epoch is finite, so this yields infinitely many actual exact refreshes.

\subsection{The value limit and the scope of the conclusion}
\label{sec:bsur-conclusion}

Every finite multiplier is at least $1-t\ge1/2$, so all individual weights stay
positive and every cache is legal. There is no fresh zero-g termination.
The proved surrogate safeguard gives
\[
 v(w+td)-v(w)\ge tG/2\ge tg/4>0
\]
at every accepted update. The auxiliary mass decreases at every update and
contracts by a factor below .99775 at each of the infinitely many fresh steps.
Thus $\omega\to0$; the invariant ratios imply $x_1,x_2\to0$ and $w_0\to1$.
By~\eqref{eq:bsur-value} and evaluation at $\theta=0$,
\[
 w_0\le v(w)\le w_0+\delta\omega+B(x_1+x_2)
 =1-.01\omega+.01(x_1+x_2)<1.
\]
This proves~\eqref{eq:bsur-claim}.

The construction is analytic; no enormous matrix or trajectory is claimed to
have been executed. It gives a second specified uncapped rule with the same
suboptimal limit, not an identity of its iterates with those of the K rule.
The conclusion concerns the unchanged G margin and this Bsur step. It does not
assert failure of every auxiliary-function method, the capped Bsur policy,
corrected-fit methods, bounded-delay policies, or uniform starts on distinct
rows. Nor does changing the coefficient make the two trajectories have the same
total work: forming $R_2$ requires additional scalar work, and their update and
refresh counts need not agree. Finally $A_w\to0$, while its scalar spectral
condition number remains one; the full information matrix becomes singular.
